% the layout is mandatory, edit only in dire circumstances

\pagestyle{empty}
\hypersetup{pageanchor=false}
\begin{center}

% top part of the layout, this actually differs between faculties

\def\ThesisTitleXmff{%
  \ifEN
    \centerline{\mbox{\includegraphics[width=166mm]{img/logo-en.pdf}}}
  \else
    \centerline{\mbox{\includegraphics[width=166mm]{img/logo-cs.pdf}}}
  \fi
  \vspace{-8mm}\vfill%
  {\bf\Large\ThesisTypeTitle}
  \vfill%
  {\LARGE\ThesisAuthor}\par
  \vspace{15mm}%
  {\LARGE\bfseries\ThesisTitle}
  \vfill%
  \Department}
\def\ThesisTitleCuniLogo#1{%
  {\large\UKName\par\medskip\par\UKFaculty }
  \vfill%
  {\bf\Large\ThesisTypeTitle}
  \vfill%
  \includegraphics[width=70mm]{#1}
  \vfill%
  {\LARGE\ThesisAuthor}\par
  \vspace{15mm}%
  {\LARGE\bfseries\ThesisTitle}
  \vfill%
  \Department\par}
\def\ThesisTitleXcuni{\ThesisTitleCuniLogo{img/uklogo.pdf}}
\def\ThesisTitleXnatur{\ThesisTitleCuniLogo{img/naturlogo.pdf}}

% choose the correct page and print it
\csname ThesisTitleX\ThesisTitleStyle\endcsname
% latex corner: X is the new @

\vfill

{
\centerline{\vbox{\halign{\hbox to 0.45\hsize{\hfil #}&\hskip 0.5em\parbox[t]{0.45\hsize}{\raggedright #}\cr
\ifEN Supervisor of the \ThesisGenitive thesis:
\else Vedoucí \ThesisGenitive práce: \fi
& \Supervisor \cr
\noalign{\vspace{2mm}}
\ifEN Study programme: \else Studijní program: \fi
& \StudyProgramme \cr
}}}}

\vfill

\ifEN Prague \else Praha \fi
\YearSubmitted

\end{center}

\newpage

% remember to sign this!
\openright
\hypersetup{pageanchor=true}
\pagestyle{plain}
\pagenumbering{roman}
\vglue 0pt plus 1fill

\ifEN
\noindent
I declare that I carried out this master thesis on my own, and only with the
cited sources, literature and other professional sources.
\else
\noindent
Prohlašuji, že jsem tuto \ThesisAccusative práci vypracoval(a) samostatně a výhradně
s~použitím citovaných pramenů, literatury a dalších odborných zdrojů.
Tato práce nebyla využita k získání jiného nebo stejného titulu.
\fi

\ifEN
\medskip\noindent
I understand that my work relates to the rights and obligations under the Act No. 121/2000 Sb., the
Copyright Act, as amended, in particular the fact that the Charles University has
the right to conclude a license agreement on the use of this work as a school work
pursuant to Section 60 subsection 1 of the Copyright Act.
\else
\medskip\noindent
Beru na~vědomí, že se na moji práci vztahují práva a povinnosti vyplývající
ze zákona č. 121/2000 Sb., autorského zákona v~platném znění, zejména skutečnost,
že Univerzita Karlova má právo na~uzavření licenční smlouvy o~užití této
práce jako školního díla podle §60 odst. 1 autorského zákona.
\fi

\vspace{10mm}


\ifEN
\hbox{\hbox to 0.5\hsize{%
In \hbox to 6em{\dotfill} date \hbox to 6em{\dotfill}
\hss}\hbox to 0.5\hsize{\dotfill\quad}}
\smallskip
\hbox{\hbox to 0.5\hsize{}\hbox to 0.5\hsize{\hfil Author's signature\hfil}}
\else
\hbox{\hbox to 0.5\hsize{%
V \hbox to 6em{\dotfill} dne \hbox to 6em{\dotfill}
\hss}\hbox to 0.5\hsize{\dotfill\quad}}
\smallskip
\hbox{\hbox to 0.5\hsize{}\hbox to 0.5\hsize{\hfil Podpis autora\hfil}}
\fi

\vspace{20mm}
\newpage

% dedication

\openright

\noindent
\Dedication

\newpage

% mandatory information page

\openright

\vbox to 0.49\vsize{\InfoPageFont
\setlength\parindent{0mm}
\setlength\parskip{5mm}

\ifEN Title: \else Název práce: \fi
\ThesisTitle

\ifEN Author: \else Autor: \fi
\ThesisAuthor

\DeptType:
\Department

\ifEN Supervisor: \else Vedoucí \ThesisGenitive práce: \fi
\Supervisor, \SupervisorsDepartment

\ifEN Abstract: \AbstractEN \else Abstrakt: \AbstractCS \fi

\ifEN Keywords: \else Klíčová slova: \fi
\Keywords

\vss}\ifEN\relax\else\nobreak\vbox to 0.49\vsize{\InfoPageFont
\setlength\parindent{0mm}
\setlength\parskip{5mm}

Title:
\ThesisTitleEN

Author:
\ThesisAuthor

\DeptTypeEN:
\DepartmentEN

Supervisor:
\Supervisor, \SupervisorsDepartmentEN

Abstract:
\AbstractEN

Keywords:
\KeywordsEN

\vss}
\fi

\newpage

\vbox to 0.49\vsize{\InfoPageFont
\setlength\parindent{0mm}
\setlength\parskip{5mm}


Název práce: Grafové neuronové sítě a hluboké posilované učení při rozvrhování prací

Autor: Maroš Bratko

Katedra: Katedra teoretické informatiky a matematické logiky

Vedoucí diplomové práce: prof. RNDr. Ing. Martin Holeňa, CSc.

Abstrakt: Prioritní rozhodovací pravidla (PDRs) jsou hladové heuristické algoritmy používané ke získaní přibližného řešení NP-těžkého úkolu organizace práce (JSP). Manuální příprava PDR vyžaduje zkušenosti v dané oblasti pro získaní dobrých výsledků. V posledních letech bylo k automatizaci přípravy PDR využito zpětnovazební učení, kde jsou PDR formulováný jako Markovův rozhodovací proces využívajíci grafovou reprezentaci JSP, ve které grafová neurónová síť (GNN) vybírá operace ke spuštění. V této prácí uvádíme přehled pěti modelů publikovaných v literatuře, kterých zdrojové kódy jsou volně dostupné na GitHubu, a experimentálňe porovnáme jejich výsledky. Naše experimentuji ukazují, že volba vstupních příznaků, napr. množství zbývající práce, významně ovlivňuje kvalitu výsledných řešení bez ohledu na architekturu GNN. To naznačuje, že výběr vstupních příznaků je pro přípravu kvalitních PDR zásadní.

Klíčová slova: organizace práce, grafové neurónové sítě, hluboké zpětnovazební učení, Markovův rozhodovací proces

\vss}

\newpage

\openright
\pagestyle{plain}
\pagenumbering{arabic}
\setcounter{page}{1}
